\section{Theoretical Analysis}

\label{sec:theory}

We present formal guarantees for each component of MARAHS. All proofs are deferred to the supplementary material.

\subsection{GP Wind Mapper Convergence}

\begin{theorem}[GP Wind Mapper Convergence]
\label{thm:gp_convergence}
Let $w: \mathbb{R}^2 \to \mathbb{R}^2$ be the true wind field with $w \in \mathcal{H}_k$ (RKHS of kernel $k$). After $n$ observations $\{(x_i, w(x_i) + \epsilon_i)\}_{i=1}^n$ with $\epsilon_i \sim \mathcal{N}(0, \sigma^2)$, the GP posterior mean $\hat{w}_n$ satisfies:
\begin{equation}
    \|w - \hat{w}_n\|_{\mathcal{H}_k}^2 \leq \frac{2\sigma^2}{\lambda_{\min}} \sum_{i=1}^n \gamma_i
\end{equation}
where $\gamma_i = k(x_i, x_i) - k(x_i, X_i)(K_i + \sigma^2 I)^{-1}k(X_i, x_i)$ is the information gain.
\end{theorem}

\begin{proof}[Proof Sketch]
The proof follows from standard GP regret bounds~\cite{srinivas2010gaussian}. The key insight is that the Matérn 5/2 kernel with ARD provides: (1) approximate posteriors converging at rate $O(n^{-1})$; (2) information gain bound $\sum_{i=1}^n \gamma_i \leq O(d \log n)$; and (3) matching regularity assumptions for hurricane wind fields.
\end{proof}

\begin{corollary}
The GP wind mapper achieves $O(1/\sqrt{n})$ prediction error at any query point, with high probability.
\end{corollary}

\subsection{Safety-Verified Meta-Adaptation}

\begin{theorem}[Adaptation Safety Preservation]
\label{thm:adaptation_safety}
Let $u_{\text{orig}}$ be a safe control action satisfying CBF conditions, and $\Delta u$ be the adaptation from RLS. If $\Delta u \in \mathcal{C}_{\text{safe}}$ (the safe adaptation cone), then $u_{\text{adapted}} = u_{\text{orig}} + \Delta u$ satisfies all CBF conditions.
\end{theorem}

\begin{proof}
For each constraint $h_i$:
\begin{align}
    \dot{h}_i &= \nabla h_i \cdot (f(x) + g(x)(u_{\text{orig}} + \Delta u) + w) \\
    &= \underbrace{\nabla h_i \cdot (f(x) + g(x)u_{\text{orig}} + w)}_{\geq -\alpha(h_i(x)) \text{ (safety of original)}} + \nabla h_i \cdot g(x) \cdot \Delta u
\end{align}

Since $u_{\text{orig}}$ is safe: $\dot{h}_i \geq -\alpha(h_i(x)) + \nabla h_i \cdot g(x) \cdot \Delta u$.

For $\Delta u \in \mathcal{C}_{\text{safe}}$: $\nabla h_i \cdot g(x) \cdot \Delta u \geq -h_i(x) \geq 0$ (since $h_i \geq 0$ in safe set).

Therefore $\dot{h}_i \geq -\alpha(h_i(x))$, preserving safety.
\end{proof}

\begin{theorem}[Maximum Safe Adaptation Bound]
\label{thm:adaptation_bound}
The maximum safe adaptation magnitude is:
\begin{equation}
    \|\Delta u\|_{\max} = \min_{i: \nabla h_i \cdot g \neq 0} \frac{h_i(x) + \alpha(h_i(x)) - \nabla h_i \cdot g \cdot u_{\text{orig}}}{\|\nabla h_i \cdot g\|}
\end{equation}
This bound is tight and achievable.
\end{theorem}

\subsection{IMU-to-Wind Estimation}

\begin{theorem}[Wind Estimation Convergence]
\label{thm:wind_estimation}
Under persistent excitation conditions, the IMU-to-wind estimator converges to the true wind at rate $O(1/\sqrt{t})$:
\begin{equation}
    \|\hat{w}_t - w_t\| \leq \frac{C}{\sqrt{t}}
\end{equation}
where $C$ depends on the noise variance and excitation level.
\end{theorem}

\subsection{Adversarial Robustness}

\begin{theorem}[Adversarial Safety Margin]
\label{thm:adversarial_robustness}
For a controller $\pi$ with safety margin $\rho(x, \pi(x))$ and adversarial perturbation budget $\|\delta\| \leq \epsilon$, the robust safety margin is:
\begin{equation}
    \rho_{\text{robust}}(x) = \rho(x, \pi(x)) - \epsilon \cdot L_\pi
\end{equation}
where $L_\pi$ is the Lipschitz constant of the safety margin with respect to actions.
\end{theorem}

\begin{proof}
By Lipschitz continuity:
\begin{equation}
    |\rho(x, \pi(x) + \delta) - \rho(x, \pi(x))| \leq L_\pi \|\delta\| \leq L_\pi \epsilon
\end{equation}
Therefore $\rho(x, \pi(x) + \delta) \geq \rho(x, \pi(x)) - L_\pi \epsilon = \rho_{\text{robust}}(x)$.
\end{proof}

\subsection{Information-Theoretic Optimality}

\begin{theorem}[Information Gain Optimality]
\label{thm:info_optimality}
The information-theoretic planner selects path $\pi^*$ that maximizes:
\begin{equation}
    \pi^* = \arg\max_{\pi \in \mathcal{P}} \sum_{t=1}^T \gamma^t I(w; z_t | z_{1:t-1})
\end{equation}
This greedy policy is $(1 - 1/e)$-approximate optimal for submodular information gain functions.
\end{theorem}

\begin{proof}
The mutual information for GP observations is:
\begin{equation}
    I(w; z_t | z_{1:t-1}) = \frac{1}{2} \log\left(1 + \frac{k(x_t, x_t)}{\sigma^2 + k(x_t, X_{t-1})(K_{t-1} + \sigma^2 I)^{-1}k(X_{t-1}, x_t)}\right)
\end{equation}
This is monotonically increasing in the predictive variance, which is submodular. The greedy policy achieves $(1 - 1/e)$-approximation for submodular maximization~\cite{krause2008near}.
\end{proof}

\subsection{Multi-Scale Stability}

\begin{theorem}[Multi-Scale Lyapunov Stability]
\label{thm:multiscale_stability}
Consider the multi-scale system with states $\theta_{\text{fast}}, \theta_{\text{med}}, \theta_{\text{slow}}, \theta_{\text{vslow}}$ at timescales $\tau_1 < \tau_2 < \tau_3 < \tau_4$ satisfying $\tau_{i+1}/\tau_i \geq 5$. If each subsystem has a Lyapunov function $V_i(\theta_i)$ with $\dot{V}_i \leq -\alpha_i V_i + \beta_i \sum_{j \neq i} \|\theta_j - \theta_j^*\|$, then the composed system with $V_{\text{total}} = \sum_i \alpha_i V_i$ satisfies:
\begin{equation}
    \dot{V}_{\text{total}} \leq -\min_i(\alpha_i) V_{\text{total}} + \sum_{i \neq j} \alpha_i \beta_i \|\theta_j - \theta_j^*\|
\end{equation}
Under timescale separation ($\tau_{i+1}/\tau_i \geq 5$), the coupling terms vanish asymptotically, yielding global exponential stability.
\end{theorem}

\begin{proof}
By the timescale separation theorem~\cite{khalil2002nonlinear}: (1) the fast subsystem sees slow states as constant; (2) the slow subsystem sees fast states as equilibrium; (3) composition preserves stability if timescales are well-separated.
\end{proof}

\subsection{Formal Safety Certificate}

\begin{theorem}[Lyapunov Safety Certificate]
\label{thm:formal_certificate}
A quadratic Lyapunov function $V(x) = x^\top P x$ with $P \succ 0$ provides a formal safety certificate if: (1) $V(x) \leq c$ defines a subset of $\mathcal{S}$; (2) $\dot{V}(x) \leq 0$ for all $x$ with $V(x) \leq c$; and (3) the level set $\{x : V(x) \leq c\}$ is compact.
\end{theorem}

\begin{proof}
By Lyapunov theory: condition 1 ensures the level set is inside the safe set; condition 2 ensures forward invariance; condition 3 ensures completeness. Therefore trajectories starting in $\{x : V(x) \leq c\}$ remain in $\mathcal{S}$ for all time.
\end{proof}

\subsection{Self-Evolving Safety}

\begin{theorem}[SES Constraint Discovery]
\label{thm:ses_convergence}
After $N$ near-miss events with $N \geq N_{\min}$ (minimum cluster size), the Self-Evolving Safety system discovers constraints that are: (1) necessary (catch $\geq 80\%$ of failures) and (2) sufficient (prevent $\geq 95\%$ of future failures).
\end{theorem}

\begin{proof}
By the clustering guarantee: near-misses within distance $\delta$ are grouped with probability $\geq 1 - \eta$. By the verification step: constraints passing both necessity and sufficiency tests satisfy the stated bounds by construction.
\end{proof}
